
% !TeX root = ../main.tex

\section{Praxisbezug und Aufwandsabschätzung}\label{sec:praxis}

Nachdem in den vorherigen Abschnitten die mathematischen
Grundlagen und der Ablauf des Schnorr-Identifikationsprotokolls
beschrieben wurden, wird im Folgenden dessen praktische
Bedeutung untersucht.

Hierzu werden zunächst der Aufwand möglicher Angriffe
und die Auswirkungen unterschiedlicher Gruppenordnungen
betrachtet. Anschließend werden konkrete Anwendungsfälle
sowie verwandte kryptografische Verfahren vorgestellt.


\subsection{Ein kleines Rechenbeispiel}

Das numerische Beispiel aus Abschnitt~\ref{sec:protokoll}
veranschaulicht die Funktionsweise des Schnorr-Protokolls
anhand der Parameter
\[
    p=23,\qquad q=11,\qquad g=2.
\]

Für den geheimen Schlüssel $x=7$, den Zufallswert $r=4$
und die Challenge $c=3$ ergibt sich das akzeptierende
Transkript
\[
    (t,c,s)=(16,3,3).
\]

Der zugehörige öffentliche Schlüssel lautet $y=13$.

Obwohl das Protokoll mathematisch korrekt ausgeführt
wird, bieten diese Parameter keine praktische
kryptografische Sicherheit.

Ein Angreifer, der lediglich die öffentlichen Werte
$p$, $q$, $g$ und $y$ kennt, kann sämtliche möglichen
privaten Schlüssel ausprobieren.

Die folgende Tabelle zeigt die Berechnung der Potenzen
des Generators $g=2$ modulo $23$.

\begin{table}[htbp]
    \centering
    \begin{tabular}{rr}
        \toprule
        Exponent $k$ & $g^k\bmod p$\\
        \midrule
        0 & 1\\
        1 & 2\\
        2 & 4\\
        3 & 8\\
        4 & 16\\
        5 & 9\\
        6 & 18\\
        7 & 13\\
        8 & 3\\
        9 & 6\\
        10 & 12\\
        \bottomrule
    \end{tabular}
    \caption{Vollständige Suche nach dem diskreten
    Logarithmus in einer kleinen zyklischen Gruppe}
    \label{tab:bruteforce-beispiel}
\end{table}

Da der öffentliche Schlüssel den Wert $y=13$ besitzt,
kann aus Tabelle~\ref{tab:bruteforce-beispiel}
unmittelbar der geheime Schlüssel
\[
    x=7
\]
abgelesen werden.

Ein Angreifer könnte sich anschließend mit diesem
Schlüssel gegenüber einem Verifier authentifizieren.

Das Beispiel zeigt, dass die mathematische Korrektheit
eines kryptografischen Protokolls allein keine
ausreichende Sicherheit garantiert.

Die Sicherheit hängt wesentlich von der Größe und
Struktur der verwendeten Gruppe sowie von der Qualität
der Implementierung ab.


\subsection{Aufwand eines Angriffs}

Für die Bewertung der Sicherheit des
Schnorr-Identifikationsprotokolls müssen unterschiedliche
Angriffsziele voneinander unterschieden werden.

Ein Angreifer kann beispielsweise versuchen,

\begin{enumerate}
    \item den langfristigen geheimen Schlüssel $x$
    aus dem öffentlichen Schlüssel $y$ zu berechnen,

    \item einen Verifier erfolgreich zu täuschen,
    ohne den geheimen Schlüssel zu kennen,

    \item Schwachstellen der Implementierung
    auszunutzen, etwa eine fehlerhafte Erzeugung
    temporärer Zufallswerte.
\end{enumerate}

Die ersten beiden Angriffsziele werden im Folgenden
näher untersucht.

Angriffe auf die Implementierung, insbesondere die
Wiederverwendung von Zufallswerten, wurden bereits in
Abschnitt~\ref{sec:protokoll} betrachtet.


\subsubsection{Schlüsselrückgewinnung durch Brute Force}

Ein grundlegender Angriff auf das Schnorr-Protokoll
besteht darin, den diskreten Logarithmus des öffentlichen
Schlüssels zu berechnen.

Gegeben sind die öffentlichen Parameter
$\G$, $q$, $g$ sowie der öffentliche Schlüssel
\[
    y=g^x.
\]

Gesucht wird der geheime Exponent
\[
    x\in\Z_q\setminus\{0\}.
\]

Bei einer vollständigen Suche, auch als Brute-Force-Angriff
bezeichnet, werden sämtliche möglichen Exponenten
systematisch überprüft.

Hierzu wird für jeden Kandidaten $k$ geprüft, ob
\[
    g^k=y
\]
gilt.

Sobald die Gleichung erfüllt ist, wurde der gesuchte
diskrete Logarithmus gefunden.

Die einzelnen Potenzen müssen dabei nicht vollständig
neu berechnet werden. Stattdessen können sie schrittweise
durch Multiplikation mit dem Generator erzeugt werden:
\[
    g^{k+1}=g^k\cdot g.
\]

Bei einem gleichverteilt gewählten geheimen Schlüssel
sind durchschnittlich etwa
\[
    \frac{q}{2}
\]
Versuche erforderlich.

Im schlechtesten Fall müssen alle $q-1$ möglichen
von null verschiedenen Exponenten überprüft werden.

Die asymptotische Laufzeit beträgt somit
\[
    \mathcal{O}(q)
\]
Gruppenoperationen.

Da bei der Suche lediglich eine konstante Anzahl
von Gruppenelementen gleichzeitig gespeichert werden
muss, beträgt der zusätzliche Speicherbedarf
\[
    \mathcal{O}(1)
\]
Gruppenelemente.

Eine vollständige Suche ist bei ausreichend großer
Gruppenordnung praktisch nicht durchführbar.
Allerdings existieren effizientere Algorithmen zur
Berechnung diskreter Logarithmen
\cite[Abschnitt 3.6.1]{menezes1996}.


\subsubsection{Baby-Step--Giant-Step}

Der Baby-Step--Giant-Step-Algorithmus ist ein
generisches Verfahren zur Berechnung diskreter
Logarithmen in zyklischen Gruppen.

Im Unterschied zur vollständigen Suche reduziert
der Algorithmus die Anzahl erforderlicher
Gruppenoperationen durch die Verwendung einer
vorberechneten Tabelle.

Hierzu wird zunächst
\[
    m=\left\lceil\sqrt{q}\right\rceil
\]
gewählt.

Jeder gesuchte Exponent $x$ mit $0\leq x<q$
kann in der Form
\[
    x=im+j
\]
mit geeigneten ganzen Zahlen $i,j$ und
$0\leq i,j<m$ dargestellt werden.

Aus
\[
    y=g^{im+j}
\]
folgt durch Umformung
\[
    y\cdot g^{-im}=g^j.
\]

Der Algorithmus berechnet zunächst die sogenannten
Baby Steps
\[
    g^j
    \qquad\text{für}\qquad
    0\leq j<m
\]
und speichert sie in einer Tabelle.

Anschließend werden die Giant Steps
\[
    y\cdot g^{-im}
\]
berechnet und mit den gespeicherten Gruppenelementen
verglichen.

Wird eine Übereinstimmung gefunden, kann der
gesuchte Exponent durch
\[
    x=im+j\bmod q
\]
bestimmt werden.

Der Algorithmus benötigt
\[
    \mathcal{O}(\sqrt{q})
\]
Gruppenoperationen sowie Speicherplatz für
\[
    \mathcal{O}(\sqrt{q})
\]
Gruppenelemente.

Bei Verwendung einer geeigneten Hashtabelle können
die gespeicherten Werte effizient nachgeschlagen
werden. Bei einer sortierten Tabelle entsteht
zusätzlicher Aufwand für die Sortierung und Suche.

Damit ist der Baby-Step--Giant-Step-Algorithmus
deutlich effizienter als eine vollständige Suche.

Sein erheblicher Speicherbedarf begrenzt jedoch
die praktische Anwendbarkeit bei großen Gruppen
\cite[Abschnitt 3.6.2]{menezes1996}.


\subsubsection{Pollard-Rho-Algorithmus}

Ein weiteres generisches Verfahren zur Berechnung
diskreter Logarithmen ist der Pollard-Rho-Algorithmus.

Dieser verwendet eine Folge pseudozufällig erzeugter
Gruppenelemente und sucht nach einer Kollision,
also nach zwei übereinstimmenden Gruppenelementen.

Die mathematische Grundlage des Verfahrens ist
mit dem Geburtstagsparadoxon verwandt.

Für eine Menge mit $q$ Elementen ist unter einer
geeigneten Zufallsannahme bereits nach ungefähr
$\sqrt{q}$ Beobachtungen mit einer Kollision
zu rechnen.

Der Pollard-Rho-Algorithmus benötigt daher im
Erwartungswert
\[
    \mathcal{O}(\sqrt{q})
\]
Gruppenoperationen.

Im Gegensatz zum Baby-Step--Giant-Step-Algorithmus
müssen jedoch nicht sämtliche bisher erzeugten
Gruppenelemente gespeichert werden.

Bei Verwendung eines geeigneten Verfahrens zur
Zyklenerkennung benötigt die klassische serielle
Variante lediglich eine konstante Anzahl von
Gruppenelementen und Skalaren:
\[
    \mathcal{O}(1).
\]

Die tatsächlich benötigte Anzahl an Operationen
ist zufallsabhängig. Ein erfolgloser Durchlauf
kann mit anderen Startwerten wiederholt werden.

Durch die Kombination aus vergleichsweise geringer
Laufzeitkomplexität und geringem Speicherbedarf
ist Pollard-Rho für die Bewertung der Sicherheit
großer zyklischer Gruppen besonders relevant
\cite[Abschnitt 3.6.3]{menezes1996}.


\subsubsection{Vergleich der Angriffsverfahren}

Tabelle~\ref{tab:dlog-algorithmen} fasst die
asymptotischen Eigenschaften der beschriebenen
generischen Angriffsverfahren zusammen.

\begin{table}[htbp]
    \centering
    \begin{tabular}{lll}
        \toprule
        Verfahren & Gruppenoperationen & Speicherbedarf\\
        \midrule
        Brute Force
            & $\mathcal{O}(q)$
            & $\mathcal{O}(1)$\\
        Baby-Step--Giant-Step
            & $\mathcal{O}(\sqrt{q})$
            & $\mathcal{O}(\sqrt{q})$\\
        Pollard-Rho
            & $\mathcal{O}(\sqrt{q})$
            & $\mathcal{O}(1)$\\
        \bottomrule
    \end{tabular}
    \caption{Vergleich generischer Algorithmen zur
    Berechnung diskreter Logarithmen.
    Die Angaben zu Pollard-Rho beziehen sich auf den
    erwarteten Aufwand der klassischen seriellen Variante.
    Der Speicherbedarf wird in Gruppenelementen und
    Skalaren angegeben. Eigene Darstellung nach
    \cite[Abschnitt 3.6]{menezes1996}.}
    \label{tab:dlog-algorithmen}
\end{table}

Aus dem Vergleich ergibt sich, dass die Sicherheit
einer zyklischen Gruppe nicht allein anhand der
Anzahl möglicher privater Schlüssel bewertet
werden darf.

Besitzt die Gruppe eine Ordnung mit ungefähr
$n$ Bit, also
\[
    q\approx 2^n,
\]
so benötigen generische Verfahren wie Pollard-Rho
ungefähr
\[
    \sqrt{q}\approx 2^{n/2}
\]
Gruppenoperationen.

Eine Gruppenordnung von ungefähr $256$ Bit
entspricht daher hinsichtlich solcher generischen
Angriffe einem Aufwand in der Größenordnung
von $2^{128}$ Gruppenoperationen.

Diese Aussage beschreibt ausschließlich die
generische Angriffskomplexität und setzt voraus,
dass keine effizienteren gruppenspezifischen
Angriffe anwendbar sind.


\subsection{Konkrete Zeitabschätzung}

Um die Größenordnung des erforderlichen
Rechenaufwands zu verdeutlichen, wird im Folgenden
eine hypothetische Zeitabschätzung vorgenommen.

Dabei wird angenommen, dass einem Angreifer eine
Rechenleistung von
\[
    R=10^{12}
\]
Gruppenoperationen pro Sekunde zur Verfügung steht.

Diese Rate ist eine bewusst gewählte Modellannahme
und kein gemessener Leistungswert realer Hardware.

Insbesondere hängt die tatsächlich erreichbare Rate
von der verwendeten Gruppe, der Größe ihrer Elemente,
der Implementierung und der eingesetzten Hardware ab.

Für einen generischen Angriff mit einem angenommenen
Aufwand von
\[
    N=\sqrt{q}
\]
Gruppenoperationen ergibt sich die geschätzte
Angriffsdauer zu
\begin{equation}\label{eq:angriffsdauer}
    T=\frac{\sqrt{q}}{R}.
\end{equation}

Konstante Faktoren, zusätzlicher Rechenaufwand und
mögliche Parallelisierungseffekte werden in dieser
vereinfachten Betrachtung nicht berücksichtigt.

\subsubsection*{Beispiel für eine Gruppenordnung von
128 Bit}

Für eine Gruppenordnung von ungefähr
\[
    q\approx 2^{128}
\]
ergibt sich ein generischer Rechenaufwand von
\[
    N\approx 2^{64}
     \approx 1{,}84\cdot 10^{19}
\]
Gruppenoperationen.

Unter der angenommenen Operationsrate beträgt die
Angriffsdauer
\begin{align*}
    T
        &\approx\frac{2^{64}}{10^{12}}\\
        &\approx 1{,}84\cdot 10^7
        \text{ Sekunden}.
\end{align*}

Dies entspricht ungefähr
\[
    T\approx 213{,}5\text{ Tagen}.
\]

Das Beispiel zeigt, dass eine Gruppenordnung von
128 Bit nicht mit einem Sicherheitsniveau von
128 Bit gleichgesetzt werden darf.


\subsubsection*{Beispiel für eine Gruppenordnung von
256 Bit}

Für eine Gruppenordnung von ungefähr
\[
    q\approx 2^{256}
\]
ergibt sich ein generischer Rechenaufwand von
\[
    N\approx 2^{128}
     \approx 3{,}40\cdot 10^{38}
\]
Gruppenoperationen.

Unter derselben angenommenen Operationsrate beträgt
die Angriffsdauer
\begin{align*}
    T
        &\approx\frac{2^{128}}{10^{12}}\\
        &\approx 3{,}40\cdot 10^{26}
        \text{ Sekunden}.
\end{align*}

Unter Verwendung von 365,25 Tagen pro Jahr ergibt
sich daraus eine Dauer von ungefähr
\[
    T\approx 1{,}08\cdot 10^{19}
    \text{ Jahren}.
\]

Diese Größenordnung liegt weit außerhalb praktisch
relevanter Zeiträume.

Tabelle~\ref{tab:angriffszeiten} vergleicht die
geschätzte Dauer für unterschiedliche Bitlängen
der Gruppenordnung.

\begin{table}[htbp]
    \centering
    \begin{tabular}{rrl}
        \toprule
        Bitlänge von $q$
            & Generischer Aufwand
            & Geschätzte Dauer\\
        \midrule
        64 Bit  & $2^{32}$  & 4,29 Millisekunden\\
        128 Bit & $2^{64}$  & 213,5 Tage\\
        160 Bit & $2^{80}$  & $3{,}83\cdot 10^4$ Jahre\\
        192 Bit & $2^{96}$  & $2{,}51\cdot 10^9$ Jahre\\
        256 Bit & $2^{128}$ & $1{,}08\cdot 10^{19}$ Jahre\\
        \bottomrule
    \end{tabular}
    \caption{Theoretische Zeitabschätzung generischer
    Angriffe bei angenommenen $10^{12}$ Gruppenoperationen
    pro Sekunde. Eigene Berechnung gemäß
    Gleichung~\eqref{eq:angriffsdauer}.}
    \label{tab:angriffszeiten}
\end{table}

Die dargestellten Werte sind keine Prognosen für
konkrete Angriffe. Sie dienen ausschließlich der
Veranschaulichung des Zusammenhangs zwischen
Gruppenordnung und generischer Angriffskomplexität.


\subsubsection{Speicherbedarf generischer Angriffe}

Neben dem zeitlichen Aufwand ist auch der
Speicherbedarf eines Angriffs von Bedeutung.

Für eine Gruppe mit einer Ordnung von ungefähr
$2^{256}$ müsste der Baby-Step--Giant-Step-Algorithmus
eine Tabelle mit ungefähr
\[
    2^{128}
\]
Gruppenelementen anlegen.

Wird beispielhaft eine Darstellung von 32 Byte
pro Gruppenelement angenommen, ergibt sich bereits
für diese Elemente ein Speicherbedarf von
\begin{align*}
    S
        &\approx 32\cdot 2^{128}\\
        &=2^{133}\text{ Byte}\\
        &\approx 1{,}09\cdot 10^{40}
        \text{ Byte}.
\end{align*}

Die Speicherung zugehöriger Exponenten sowie
zusätzlicher Tabellenstrukturen würde den Bedarf
weiter erhöhen.

Die angenommene Elementgröße dient lediglich der
Illustration. Bei einer multiplikativen Untergruppe
eines großen endlichen Körpers können Gruppenelemente
wesentlich mehr Speicherplatz benötigen.

Der Pollard-Rho-Algorithmus besitzt hingegen einen
wesentlich geringeren Speicherbedarf und ist
deshalb für generische Sicherheitsabschätzungen
besonders relevant.


\subsubsection{Spezialisierte Angriffe auf endliche Körper}

Die bisher betrachteten Algorithmen sind generisch,
da sie keine besonderen algebraischen Eigenschaften
der verwendeten Gruppe ausnutzen.

Für bestimmte Gruppen existieren jedoch
spezialisierte Verfahren, die effizienter als
generische Algorithmen sein können.

Bei multiplikativen Gruppen endlicher Körper zählen
hierzu insbesondere Index-Calculus-Verfahren und
darauf aufbauende Varianten des Number Field Sieve.

Diese Verfahren nutzen die konkrete mathematische
Darstellung der Gruppenelemente aus.

Für große Primkörper ist der asymptotische Aufwand
des allgemeinen Number Field Sieve zur Berechnung
diskreter Logarithmen subexponentiell in der
Bitlänge des Körpermodulus.

Daher hängt die Sicherheit einer Untergruppe von
$\Z_p^*$ nicht nur von ihrer Gruppenordnung $q$,
sondern auch wesentlich von der Größe und Struktur
des Primzahlmodulus $p$ ab
\cite[Kapitel 15]{galbraith2012}.

Eine große Gruppenordnung allein garantiert
in diesem Fall kein entsprechendes Sicherheitsniveau.

Bei geeigneten elliptischen Kurven sind hingegen
keine klassischen subexponentiellen Algorithmen
zur Lösung des allgemeinen elliptischen diskreten
Logarithmusproblems bekannt.

Für solche Kurven stellen generische Angriffe
wie Pollard-Rho daher eine wesentliche Grundlage
der Sicherheitsabschätzung dar.

Diese Aussage gilt nicht uneingeschränkt für
jede elliptische Kurve. Bestimmte spezielle Kurven
können zusätzliche mathematische Schwächen
aufweisen.

Für die praktische Auswahl kryptografischer
Parameter müssen deshalb sowohl die Gruppenordnung
als auch die Struktur der zugrunde liegenden
Gruppe berücksichtigt werden.


\subsection{Erfolgreiche Täuschung ohne Schlüsselkenntnis}

Neben der vollständigen Rückgewinnung des geheimen
Schlüssels besteht ein weiteres Angriffsziel darin,
eine einzelne Protokollausführung erfolgreich
abzuschließen, ohne den privaten Schlüssel zu kennen.

Wie in Abschnitt~\ref{sec:protokoll} gezeigt wurde,
kann ein Teilnehmer ein akzeptierendes Transkript
für eine im Voraus festgelegte Challenge erzeugen.

Hierzu wählt er eine Challenge $c'$ und eine
Response $s'$ und berechnet
\[
    t'=g^{s'}y^{-c'}.
\]

Sendet er dieses Commitment an einen ehrlichen
Verifier, so kann er die vorbereitete Response
verwenden, falls die tatsächliche Challenge
mit der zuvor gewählten Challenge übereinstimmt.

Für eine unabhängige und gleichverteilt aus
$\Z_q$ gewählte Challenge gilt
\begin{equation}\label{eq:taeuschungswahrscheinlichkeit}
    \Pr[c=c']=\frac{1}{q}.
\end{equation}

Dies ist die Erfolgswahrscheinlichkeit dieser
konkreten Strategie und entspricht dem in
Abschnitt~\ref{sec:protokoll} beschriebenen
Wissensfehler.

Sie darf nicht als allgemeine obere Schranke
für sämtliche denkbaren Angriffe interpretiert
werden.

\paragraph{Beispiel mit kleinen Parametern.}

Für die verwendete Lernumgebung mit $q=11$
beträgt die Erfolgswahrscheinlichkeit
\[
    \Pr[c=c']
        =\frac{1}{11}
        \approx 9{,}09\,\%.
\]

Ein Angreifer hätte mit dieser Strategie
somit bereits bei einer einzelnen Protokollrunde
eine vergleichsweise hohe Erfolgswahrscheinlichkeit.

\paragraph{Beispiel mit realistischen Größenordnungen.}

Bei einer Challenge aus einem Wertebereich mit
$2^{128}$ gleichwahrscheinlichen Möglichkeiten
beträgt die Erfolgswahrscheinlichkeit des
einfachen Ratens
\[
    \Pr[c=c']=2^{-128}
    \approx 2{,}94\cdot 10^{-39}.
\]

Werden mehrere unabhängige Authentifizierungsversuche
zugelassen, steigt die Wahrscheinlichkeit,
mindestens einen erfolgreichen Versuch zu erzielen.

Für $k$ unabhängige Versuche mit einer jeweiligen
Erfolgswahrscheinlichkeit von $1/q$ ergibt sich
\[
    P_k=1-\left(1-\frac{1}{q}\right)^k.
\]

Für $k\ll q$ kann diese Wahrscheinlichkeit durch
\[
    P_k\approx\frac{k}{q}
\]
angenähert werden.

In praktischen Anwendungen können zusätzliche
Maßnahmen, beispielsweise die Begrenzung der
Anzahl von Authentifizierungsversuchen, das Risiko
solcher Online-Angriffe weiter reduzieren.


\subsection{Rechenaufwand des legitimen Protokolls}

Neben dem Aufwand möglicher Angriffe ist auch die
Effizienz des Protokolls für legitime Teilnehmer
von Bedeutung.

Das Schnorr-Identifikationsprotokoll ist
vergleichsweise effizient, da es nur wenige
aufwendige Gruppenoperationen benötigt.

Der Prover führt pro Protokollausführung im
Wesentlichen folgende Berechnungen durch:
\[
    t=g^r
\]
und
\[
    s=(r+cx)\bmod q.
\]

Die Berechnung des Commitments erfordert eine
Gruppenexponentiation.

Die Berechnung der Response besteht hingegen
lediglich aus einer modularen Multiplikation
und einer modularen Addition.

Der Verifier überprüft anschließend
\[
    g^s=t\,y^c.
\]

Eine direkte Implementierung benötigt hierfür
zwei Gruppenexponentiationen und eine
Gruppenmultiplikation.

Durch Verfahren zur simultanen Exponentiation
kann der Rechenaufwand der Verifikation
reduziert werden.

Tabelle~\ref{tab:legitimer-aufwand} fasst die
wesentlichen Berechnungen zusammen.

\begin{table}[htbp]
    \centering
    \begin{tabular}{lll}
        \toprule
        Partei & Berechnung & Aufwand\\
        \midrule
        Prover
            & $t=g^r$
            & Eine Exponentiation\\
        Prover
            & $s=(r+cx)\bmod q$
            & Skalare Arithmetik\\
        Verifier
            & $g^s\stackrel{?}{=}t\,y^c$
            & Zwei Exponentiationen\\
        \bottomrule
    \end{tabular}
    \caption{Wesentliche Berechnungen einer regulären
    Schnorr-Identifikation ohne zusätzliche
    Parameter- und Schlüsselvalidierung.
    Eigene Darstellung.}
    \label{tab:legitimer-aufwand}
\end{table}

Die Anzahl benötigter Gruppenoperationen für
eine einzelne Exponentiation wächst bei Verwendung
geeigneter Algorithmen, beispielsweise Square-and-Multiply,
nur logarithmisch mit der Größe des Exponenten.

Bei einem Exponenten aus $\Z_q$ ergibt sich
ein Aufwand von
\[
    \mathcal{O}(\log q)
\]
Gruppenoperationen.

Damit besteht ein erheblicher Unterschied
zwischen dem Aufwand der legitimen Teilnehmer
und dem Aufwand eines generischen Angreifers:
\[
    \mathcal{O}(\log q)
    \qquad\text{gegenüber}\qquad
    \mathcal{O}(\sqrt{q}).
\]

Diese Asymmetrie ist eine wesentliche Grundlage
der praktischen Einsetzbarkeit des Verfahrens.

Weitere Ausführungen zur Effizienz der verwandten
nichtinteraktiven Variante finden sich in
\cite[Abschnitte 2.4 und 3.4]{hao2017}.


\subsection{Anwendungsfälle und verwandte Verfahren}

Das Schnorr-Identifikationsprotokoll stellt ein
grundlegendes kryptografisches Verfahren dar,
das sowohl zur direkten Identifikation als
auch als Baustein komplexerer Verfahren
eingesetzt werden kann.

Dabei muss zwischen dem ursprünglichen interaktiven
Protokoll und den daraus abgeleiteten Verfahren
unterschieden werden.


\subsubsection{Interaktive Schlüsselauthentifizierung}

Ein grundlegender Anwendungsfall des
Schnorr-Identifikationsprotokolls ist die
Authentifizierung eines Teilnehmers anhand
eines asymmetrischen Schlüsselpaares.

Bei einer Registrierung erzeugt der Teilnehmer
einen privaten Schlüssel $x$ und hinterlegt
den zugehörigen öffentlichen Schlüssel
\[
    y=g^x.
\]

Bei einer späteren Anmeldung weist der Teilnehmer
mithilfe des interaktiven Schnorr-Protokolls
nach, dass er den zum registrierten öffentlichen
Schlüssel gehörenden privaten Schlüssel kennt.

Der private Schlüssel wird dabei nicht an
den authentifizierenden Server übertragen.

Dieser Ansatz unterscheidet sich von einer
klassischen passwortbasierten Authentifizierung,
bei der ein gemeinsames Geheimnis beziehungsweise
eine daraus abgeleitete Information zur
Überprüfung verwendet wird.

Bei einer schlüsselbasierten Authentifizierung
muss der Server hingegen lediglich den
öffentlichen Schlüssel zur Überprüfung
des mathematischen Nachweises speichern.

Die Offenlegung des öffentlichen Schlüssels
ermöglicht unter der Annahme der Schwierigkeit
des diskreten Logarithmusproblems keine
effiziente Berechnung des privaten Schlüssels.

Allerdings sind für eine sichere praktische
Authentifizierung zusätzliche Mechanismen
erforderlich.

Dazu zählen insbesondere eine vertrauenswürdige
Registrierung des öffentlichen Schlüssels, der
Schutz des privaten Schlüssels und die Bindung
der Authentifizierung an die vorgesehene
Gegenstelle.

Die hier beschriebene Anwendung stellt einen
konzeptionellen Anwendungsfall des interaktiven
Protokolls dar und ist nicht mit einem bestimmten
standardisierten Anmeldedienst gleichzusetzen.


\subsubsection{Nichtinteraktive Wissensnachweise}

Eine wesentliche Erweiterung des Schnorr-Protokolls
ist der nichtinteraktive Zero-Knowledge-Nachweis.

Im Gegensatz zur ursprünglichen interaktiven
Variante muss der Verifier keine zufällige
Challenge an den Prover übertragen.

Stattdessen wird die Challenge mittels einer
kryptografischen Hashfunktion aus dem Commitment,
dem öffentlichen Schlüssel und weiteren
Kontextinformationen abgeleitet.

Diese Transformation geht auf das
Fiat-Shamir-Verfahren zurück und wurde bereits
in Abschnitt~\ref{sec:protokoll} vorgestellt.

RFC~8235 beschreibt eine konkrete nichtinteraktive
Variante des Schnorr-Protokolls sowohl für
multiplikative Untergruppen endlicher Körper
als auch für elliptische Kurven.

Das Dokument ist als informatorischer RFC
veröffentlicht und beschreibt unter anderem
Anwendungsfälle sowie sicherheitsrelevante
Implementierungsvoraussetzungen.

Nichtinteraktive Wissensnachweise können als
Baustein kryptografischer Protokolle verwendet
werden, um die Kenntnis eines diskreten
Logarithmus nachzuweisen, ohne den zugrunde
liegenden Exponenten offenzulegen
\cite[Abschnitte 2, 3 und 5]{hao2017}.


\subsubsection{Passwortauthentifizierter Schlüsselaustausch}

Ein konkreter Anwendungsfall des nichtinteraktiven
Schnorr-Wissensnachweises ist das J-PAKE-Protokoll
(\emph{Password-Authenticated Key Exchange by Juggling}).

J-PAKE ermöglicht es zwei Parteien, auf Grundlage
eines gemeinsam bekannten Passworts einen
kryptografischen Sitzungsschlüssel auszuhandeln.

Hierzu tauschen die Teilnehmer unter anderem
Gruppenelemente aus, deren zugehörige Exponenten
sie mithilfe von Schnorr-Wissensnachweisen
nachweisen.

Diese Nachweise dienen dazu, die korrekte
Erzeugung bestimmter Protokollwerte zu überprüfen,
ohne die jeweiligen geheimen Exponenten zu
übertragen.

Das Verfahren wird in RFC~8236 beschrieben
\cite[Abschnitte 1 und 2]{hao2017jpake}.

Der Anwendungsfall verdeutlicht, dass Schnorr-Nachweise
auch innerhalb komplexerer kryptografischer
Protokolle eingesetzt werden können.

Dabei handelt es sich nicht um eine unmittelbare
Ausführung des ursprünglichen interaktiven
Identifikationsprotokolls, sondern um die
Integration seiner nichtinteraktiven Variante
in ein umfassenderes Schlüsselaustauschverfahren.


\subsubsection{Nachweis des Besitzes eines privaten
Schlüssels}

Ein weiterer möglicher Anwendungsfall ist der
sogenannte \emph{Proof of Possession} (PoP)
bei der Registrierung eines öffentlichen
Schlüssels.

Beispielsweise kann eine Zertifizierungsstelle
vor der Ausstellung eines digitalen Zertifikats
einen Nachweis verlangen, dass der Antragsteller
tatsächlich über den zum eingereichten öffentlichen
Schlüssel gehörenden privaten Schlüssel verfügt.

Ein Schnorr-Wissensnachweis kann für einen
solchen Nachweis verwendet werden, sofern er
eindeutig an den jeweiligen Registrierungsvorgang
gebunden ist.

Hierzu können beispielsweise Informationen
über die Zertifizierungsstelle, den Antragsteller
oder den konkreten Vorgang in die Berechnung
der nichtinteraktiven Challenge einbezogen
werden.

RFC~8235 nennt diesen Anwendungsfall ausdrücklich
als mögliche Verwendung des Schnorr-Nachweises
\cite[Abschnitt 5]{hao2017}.

Dadurch wird nicht automatisch die tatsächliche
Identität des Antragstellers nachgewiesen.

Der Nachweis bezieht sich zunächst ausschließlich
auf die Kenntnis des zum öffentlichen Schlüssel
gehörenden privaten Schlüssels.

Die Überprüfung der Identität stellt einen
zusätzlichen organisatorischen und technischen
Prozess dar.


\subsubsection{Schnorr-Signaturen}

Eine weitere wichtige Entwicklung auf Grundlage
des Schnorr-Identifikationsprotokolls sind
digitale Schnorr-Signaturen.

Bei digitalen Signaturen wird der Nachweis
mit einer konkreten Nachricht verknüpft.

In einer vereinfachten Darstellung wird die
Challenge aus dem Commitment, dem öffentlichen
Schlüssel und der zu signierenden Nachricht
$m$ abgeleitet:
\[
    c=H(t,y,m)\bmod q.
\]

Anschließend wird die Response wie beim
interaktiven Protokoll berechnet:
\[
    s=(r+cx)\bmod q.
\]

Eine mögliche abstrakte Darstellung der
Signatur ist
\[
    \sigma=(t,s).
\]

Der Verifier berechnet die Challenge aus
den öffentlichen Informationen erneut und
überprüft die zugehörige Verifikationsgleichung.

Diese Darstellung verdeutlicht lediglich
den mathematischen Zusammenhang zwischen
Schnorr-Identifikation und Schnorr-Signaturen.

Konkrete Signaturstandards definieren zusätzliche
Vorgaben für die Kodierung der Werte,
die Erzeugung temporärer Skalare und die
Verwendung der Hashfunktion.

Ein praktisches Beispiel für den Einsatz
von Schnorr-Signaturen ist Bitcoin.

Mit dem Bitcoin Improvement Proposal 340
(BIP~340) wurde ein Schnorr-Signaturverfahren
auf Grundlage der elliptischen Kurve
secp256k1 spezifiziert
\cite{wuille2020bip340}.

Schnorr-Signaturen werden dort im Rahmen von
Taproot zur Autorisierung entsprechender
Transaktionen eingesetzt.

Dieser Anwendungsfall zeigt die praktische
Bedeutung der mathematischen Prinzipien des
Schnorr-Protokolls, obwohl nicht das
ursprüngliche interaktive Identifikationsverfahren
ausgeführt wird.


\subsubsection{Abgrenzung zu modernen Passkeys}

Auch moderne passwortlose Authentifizierungsverfahren,
wie beispielsweise Passkeys, verwenden asymmetrische
Schlüsselpaare und kryptografische Nachweise
über den Besitz eines privaten Schlüssels.

Dadurch besteht eine konzeptionelle Ähnlichkeit
zum Schnorr-Identifikationsprotokoll.

Passkeys sind jedoch nicht mit dem
Schnorr-Protokoll gleichzusetzen.

Sie beruhen auf den durch FIDO2 und WebAuthn
definierten Authentifizierungsmechanismen und
können unterschiedliche digitale
Signaturverfahren verwenden.

Darüber hinaus werden die Authentifizierungsdaten
an den jeweiligen Webdienst und den
Authentifizierungsvorgang gebunden.

Die Gemeinsamkeit besteht somit vor allem
in der Verwendung asymmetrischer Kryptografie
zur Authentifizierung, ohne den privaten
Schlüssel an den Server übertragen zu müssen.

Eine direkte Verwendung des interaktiven
Schnorr-Protokolls durch Passkeys kann daraus
nicht abgeleitet werden.


\subsection{Zusammenfassung}

Die Untersuchung des Angriffsaufwands zeigt,
dass die Sicherheit des Schnorr-Protokolls
wesentlich von den gewählten kryptografischen
Parametern abhängt.

Während eine vollständige Suche nach dem
geheimen Schlüssel einen Aufwand von
$\mathcal{O}(q)$ Gruppenoperationen besitzt,
reduzieren generische Algorithmen wie
Baby-Step--Giant-Step und Pollard-Rho diesen
Aufwand auf $\mathcal{O}(\sqrt{q})$.

Für bestimmte mathematische Gruppen können
darüber hinaus spezialisierte Angriffe
existieren.

Eine sichere Parameterauswahl muss daher
sowohl die Gruppenordnung als auch die
Struktur der zugrunde liegenden Gruppe
berücksichtigen.

Gleichzeitig benötigen die legitimen Teilnehmer
des Protokolls vergleichsweise wenige
Gruppenoperationen.

Diese Kombination aus effizienter regulärer
Protokollausführung und hohem angenommenem
Angriffsaufwand macht die zugrunde liegenden
mathematischen Prinzipien für praktische
kryptografische Anwendungen interessant.

Die betrachteten Anwendungsfälle verdeutlichen
zudem, dass das Schnorr-Protokoll nicht nur
als eigenständiges Identifikationsverfahren,
sondern insbesondere als Grundlage für
nichtinteraktive Wissensnachweise, komplexere
Schlüsselaustauschprotokolle und digitale
Signaturen von Bedeutung ist.
